% !TeX root = ../main.tex
% =====================================================================
% Method section — ReCAP (backprop formulation), v2 (review revision).
%   * Written against the \soap macro; labels sec:method:recap and
%     eq:recap (renamed from sec:method:crr / eq:crr — grep old names).
%   * Review edits: mean pooling only (no \phi / last-token option);
%     mean-centering removed; orientation abstracted to a map \rho;
%     set-size normalizations |.| replaced by \operatorname*{mean}
%     (needs amsmath); spectrum-structure and attention-asymmetry
%     explanations commented out, not deleted; Properties paragraph
%     commented out; contamination paragraph simplified.
% =====================================================================

% \note{add the ablation for context cutoff in appendix. context ablation: 1024, 2048, 4096, 8192 (what we use), 16384, on WW-AG and WW-HC only.}
% \textcolor{blue}{Please don't mention the details here. add ablations in appendix.}


% Method figure updated 2026-09-01: artifacts/method_figure/soap_framework.{pptx,pdf}
% -> manuscript/assets/soap_framework.pdf (two-stage schematic; matrices, scores
% and weights illustrative). Supersedes assets/soap_overview.pdf (2026-08-25).

% Let $\mathcal{T}_t$ denote the token positions corresponding to $s_t$ in
% $z_t$, and let $h_p^{(l)}\in\mathbb{R}^d$ be the hidden representation at
% layer $l$ and token position $p$.
% We obtain the $d$-dimensional step representation by mean-pooling over the tokens belonging
% to the current step:
% $
%     v_t
%     =
%     \frac{1}{|\mathcal{T}_t|}
%     \sum_{p\in\mathcal{T}_t}
%     h_p^{(l)}
%     \in\mathbb{R}^d.
% $



% \note{should clarify? in practice, we trim the context by removing early steps until $z_{t}$ fits into the context limit}\textcolor{blue}{put in the appendix}


\section{\soap: Spectral Scoring with Attention-Guided Propagation}
\label{sec:method}


We introduce \soap, a framework for localizing decisive errors from
step-unlabeled failure trajectories.
At a high level, \soap consists of two stages.
\textit{First}, we construct a spectral reference space from step
representations extracted by $\proxy$ on $\mathcal{D}_{\mathrm{fail}}$ and
derive a base error score for each step
(\S\ref{sec:method:base}).
\textit{Second}, we estimate soft inter-step dependencies from the attention
patterns of $\proxy$ and propagate error evidence from anomalous downstream
steps back to their likely upstream sources
(\S\ref{sec:method:weights}).
The resulting score $\Stilde(s_t)$ measures how likely step $s_t$ is to have
introduced the failure.
Figure~\ref{fig:overview} summarizes the framework.



\subsection{Spectral Step Scoring from Unlabeled Failure Trajectories}
\label{sec:method:base}


Although every trajectory in $\mathcal{D}_{\mathrm{fail}}$ ends in failure,
most individual steps remain locally valid, while failure-related steps
constitute only a small fraction of each trajectory~\citep{whoandwhen}.
Consequently, the representation space induced by $\proxy$ is dominated by
recurring patterns of ordinary agent behavior, with failure-related steps
forming a sparse, unlabeled component.
We exploit this imbalance by spectrally decomposing the  reference
representations and measuring how strongly each step aligns with selected
spectral directions.


For a trajectory $\traj=(s_1,\ldots,s_T)$, we encode each step using the task
description $x$ and selected trajectory context associated with that step.
Let
$
    v_t\in\mathbb{R}^{d}
$
denote the resulting representation extracted by $\proxy$ for step $s_t$ (cf. Appendix~\ref{app:implementation}).



\textbf{Spectral reference space.}
We apply the representation extraction procedure to every step in
$\mathcal{D}_{\mathrm{fail}}$, obtaining
$v_t^{(n)}\in\mathbb{R}^{d}$ for the $t$-th step of the $n$-th reference
trajectory.
Let
$
    M_{\mathrm{ref}}
    =
    \sum_{n=1}^{N}T_n
$
denote the total number of reference steps.
Stacking these representations row-wise gives
$R\in\mathbb{R}^{M_{\mathrm{ref}}\times d}$.
Let $r=\min\{M_{\mathrm{ref}},d\}$.
We compute the compact singular value decomposition
\begin{equation}
\small
    R
    =
    U\Sigma V^\top,
    \qquad
    V\in\mathbb{R}^{d\times r}.
    \label{eq:reference-svd}
\end{equation}
The columns of $V$ define orthogonal directions in the reference
representation space, ordered by non-increasing singular value.


Because ordinary steps dominate $R$, the leading singular directions may
primarily encode frequent shared structure, such as response style,
trajectory format, and common tool-use patterns.
Conversely, directions at the extreme tail may capture irrelevant
variation.
We therefore do not restrict scoring to either end of the spectrum.
Instead, we select a contiguous spectral band
\begin{equation}
\small
    \mathcal{C}
    =
    \{c_{\mathrm{begin}},\ldots,c_{\mathrm{end}}\},
    \qquad
    1
    \leq
    c_{\mathrm{begin}}
    \leq
    c_{\mathrm{end}}
    \leq
    r,
    \label{eq:spectral-band}
\end{equation}
whose boundaries are selected as described in
\S\ref{app:implementation}.

% \textbf{Spectral reference space.}
% We apply the same representation extraction procedure to every step in the
% reference collection $\mathcal{D}_{\mathrm{fail}}$.
% Let
% $
%     M=\sum_{n=1}^{N}T_n
% $
% be the total number of reference steps.
% We stack their representations into $R=[(v_t^{(n)})^\top]_{n\in[N],\,t\in[T_n]}\in\mathbb{R}^{M\times d}$ where $v_t^{(n)}$ is representation of the $t$-th step in the $n$-th sequence or trajectory. We then compute its  singular value decomposition
% \begin{equation}
%     R = U\Sigma V^\top.
%     \label{eq:reference-svd}
% \end{equation}
% The columns of $V$ describe dominant orthogonal directions in the
% LLM representation space, ordered by decreasing singular value. 

\begin{wrapfigure}{r}{0.35\linewidth}
\vspace{-0.9\baselineskip}
\centering
\includegraphics[width=\linewidth]{{assets/qwen3.5-9b_arc_seeds-6-7-8_seed-7}.pdf}
\vspace{-0.8cm}
\caption{\small \textbf{Base-score distributions} on a \textsc{Correct-Error} test
split~\citep{correct}.
Decisive-error steps exhibit lower projection energy than ordinary steps.}
\label{fig:score-dist}
\vspace{-\baselineskip}
\end{wrapfigure}

% Importantly, the reference set $R$ is dominated by ordinary steps because
% each failure trajectory typically contains \textit{only a small number of
% decisive errors}.
% Consequently, the leading singular directions may primarily encode
% frequent shared structure, such as  response style,
% trajectory format, and common tool-use patterns.
% Conversely, components at the extreme tail may contain irrelevant variation.
% We therefore do not assume that the most informative signal must lie in
% the leading components.
% Instead, we consider a contiguous spectral band to represent the decisive-error information
% $
%     \mathcal{C}
%     =
%     \{c_{\mathrm{begin}},\ldots,c_{\mathrm{end}}\},
% $ whose boundaries are selected as described in
% % \S\ref{sec:experiments}.
% \S\ref{app:implementation}.
% \textcolor{red}{which precise section?} \note{cited.}

\textbf{Spectral base score.}
For a step representation $v_t$, we compute its average squared projection
onto the selected spectral band and invert this quantity:
\begin{equation}
\small
    q_{\mathcal{C}}(v_t)
    =
    \frac{1}{|\mathcal{C}|}
    \sum_{c\in\mathcal{C}}
    \left\langle v_t,V_{:,c}\right\rangle^2,\quad
    \Sbase(s_t)
    =
    \frac{1}{q_{\mathcal{C}}(v_t)+\epsilon},
\label{eq:base-score}
\end{equation}
where $\epsilon>0$ ensures numerical stability.
Because the reference spectrum is learned primarily from ordinary steps,
decisive-error steps tend to exhibit lower projection energy within the
selected band
(Figure~\ref{fig:score-dist}).
The reciprocal therefore assigns larger base scores to steps exhibiting
stronger decisive-error evidence.
Constructing $R$ and its spectral decomposition \textit{uses no decisive-error or
responsible-agent annotations.}


% \textbf{Spectral base score.}
% For a step representation $v_t$, we first compute its average squared
% projection onto the selected spectral band and then invert this quantity:
% \begin{equation}
% q_{\mathcal C}(v_t)
% =
% \frac{1}{|\mathcal C|}
% \sum_{c\in\mathcal C}
% \left\langle v_t,V_{:,c}\right\rangle^2,\quad
% \Sbase(s_t)
% =
% \frac{1}{q_{\mathcal C}(v_t)+\epsilon},
% \label{eq:base-score}
% \end{equation}
% where $\epsilon>0$ ensures numerical stability.
% Because ordinary steps dominate the reference corpus, recurring patterns of
% ordinary behavior are more strongly represented in its spectral structure.
% Decisive-error steps therefore tend to have lower projection energy within
% the selected band (Figure~\ref{fig:score-dist}), yielding larger base scores.
% The reference matrix $R$ is constructed exclusively from
% $\mathcal D_{\mathrm{fail}}$, \textit{without decisive-error or responsible-agent
% annotations.}




\subsection{Dependency-aware Rescoring}
\label{sec:method:weights}

\textbf{Attention-guided step dependencies.} The base score in Eq.~\ref{eq:base-score} evaluates each step independently
and therefore cannot distinguish an error source from its downstream
consequences.
Once a decisive error occurs, subsequent steps may inherit the corrupted
context and exhibit stronger failure signals than the original error.
An earlier step is therefore a more plausible root cause when anomalous
downstream steps depend strongly on it.
We capture this structure through a soft dependency graph over the target
trajectory.

For each step $s_t$, let
$
    \mathcal{P}_t
    \subseteq
    \{1,\ldots,t-1\}
$
denote the index set of eligible predecessor steps retained by our context
construction.
For each $i\in\mathcal{P}_t$, we aggregate the post-softmax attention routed
from tokens in $s_t$ to tokens in $s_i$, yielding a non-negative attention
mass $m_{i,t}$.
We normalize these masses over the retained predecessors:
\begin{equation}
\small
    w_{i,t}
    =
    \begin{cases}
        \displaystyle
        \frac{m_{i,t}}
             {\sum_{j\in\mathcal{P}_t}m_{j,t}},
        & i\in\mathcal{P}_t,\\[1.2ex]
        0,
        & \text{otherwise}.
    \end{cases}
    \label{eq:dependency-weight}
\end{equation}
Thus,
$
    \sum_{i\in\mathcal{P}_t}w_{i,t}=1
$
whenever $\mathcal{P}_t$ is nonempty, and $w_{i,t}$ represents the relative
dependency of the later step $s_t$ on the earlier step $s_i$.
The construction of $\mathcal{P}_t$ and the attention aggregation procedure
are detailed in Appendix~\ref{app:implementation}.

These weights provide model-internal evidence about which earlier steps
contribute to the representation of a later step.
Because they are obtained directly from the frozen LLM model, their
construction requires neither repeated generative prompting, which
introduces substantial inference cost~\citep{whoandwhen,oat}, nor explicit
dependency reconstruction and traversal through LLM-based graph
reasoning~\citep{falat,chief,bakish2026adaptive}, nor task-specific
post-training on attributed failure trajectories
~\citep{agentracer,graphtracer}.






% The base score in Eqn.~\ref{eq:base-score} evaluates each step independently
% and therefore cannot distinguish an error source from its downstream
% consequences.
% Once a decisive error occurs, later steps may inherit the corrupted
% context and exhibit stronger failure signals than the original error.
% An earlier step is therefore more plausible as the root cause when
% subsequent anomalous steps depend strongly on it.
% To capture this structure, we construct a soft dependency graph over
% the target trajectory.



% For every pair $(i,t)$ with $i<t$, we define a non-negative coefficient
% $w_{i,t}$ that measures how strongly the later step $s_t$ depends on the
% earlier step $s_i$.
% We derive this coefficient from the post-softmax attention produced by
% $\proxy$. Specifically, we aggregate the attention routed from tokens in $s_t$ to
% tokens in $s_i$, yielding
% an attention mass $m_{i,t}$.
% We then normalize over all predecessor steps:
% \begin{equation}
%     w_{i,t}
%     =
%     \frac{m_{i,t}}
%          {\sum_{j<t}m_{j,t}},
%     \qquad i<t.
%     \label{eq:dependency-weight}
% \end{equation}
% Thus, $\sum_{i<t}w_{i,t}=1$, and $w_{i,t}$ represents the relative
% dependency of step $s_t$ on predecessor $s_i$. The resulting weights provide soft evidence about which preceding
% steps contribute to the LLM model's representation of a later step. They require neither repeated LLM prompting, which introduces
% substantial inference cost~\citep{whoandwhen,oat}, nor explicit dependency
% reconstruction and traversal through LLM-based graph reasoning~\citep{falat,chief,bakish2026adaptive}, nor task-specific post-training on attributed
% failure trajectories~\citep{agentracer,graphtracer}. 
% \note{in practice, we don't use all dependencies, indicated by the hyperparameter window size $w$. refer this details in the appendix.}






\textbf{Rescoring.} We treat the error signal of a downstream step as additional evidence for
the earlier steps on which it depends.
For each step $s_i$, we compute the dependency-weighted average of its downstream base
scores:
\begin{equation}
\small
    B_i
    =
        \frac{
            \sum_{t>i}
            w_{i,t}\,\Sbase(s_t)
        }{
           \sum_{t>i}w_{i,t}
        },
    \label{eq:backward-evidence}
\end{equation}
In particular, $B_T=0$ because the final step has no downstream dependents. We combine the local error signal with the downstream evidence attributed to
the step:
\begin{equation}
\small
    \Stilde(s_i)
    =
    \Sbase(s_i)
    +
    \gamma B_i,
    \qquad
    \gamma\in[0,1].
    \label{eq:recap}
\end{equation}
Here, $\gamma$ controls the strength of dependency-aware rescoring, with
$\gamma=0$ recovering the spectral base score. The normalization in Eq.~\ref{eq:backward-evidence} prevents an early step
from receiving a large correction merely because it has many downstream
dependents.
Instead, downstream steps with larger dependency weights exert greater
influence on its correction.
Thus, an earlier step is promoted when the downstream steps that depend
strongly on it exhibit large base error scores.

Finally, we predict the decisive-error step and corresponding responsible
agent as
\begin{equation}
\small
    \hat{t}
    =
    \arg\max_{i\in\{1,\ldots,T\}}
    \Stilde(s_i),
    \qquad
    \hat{a}
    =
    a_{\hat{t}}.
    \label{eq:prediction}
\end{equation}



% Given the dependency weights, we treat the error signal of a downstream
% step as additional evidence for the earlier steps on which it depends.
% For each step $s_i$, we compute the dependency-weighted average of the
% base scores of its downstream dependents:
% \begin{equation}
%     B_i
%     =
%     \frac{
%         \sum_{t>i}
%         w_{i,t}\,\Sbase(s_t)
%     }{
%         \sum_{t>i}w_{i,t}
%     },
%     \qquad
%     B_T=0.
%     \label{eq:backward-evidence}
% \end{equation}
% We then combine the step's local error signal with the downstream evidence attributed to it:
% \begin{equation}
%     \Stilde(s_i)
%     =
%     \Sbase(s_i)
%     +
%     \gamma B_i,
%     \qquad
%     \gamma\in[0,1].
%     \label{eq:recap}
% \end{equation}
% Here, $\gamma$ controls the strength of dependency-aware rescoring, and
% $\gamma=0$ recovers the base score.

% The normalization in Eqn.~\ref{eq:backward-evidence} produces a
% dependency-weighted average, which prevents an early or frequently attended step from receiving a
% large correction merely because it has many downstream steps.
% Instead, a step is promoted when the later steps that depend on it
% consistently exhibit strong failure signals.
% Conversely, an anomalous downstream step contributes little evidence to
% a predecessor on which it places little dependency weight.


% Finally, we predict the decisive-error step and the corresponding
% responsible agent as
% \begin{equation}
%     \hat{t}
%     =
%     \arg\max_{i\in\{1,\ldots,T\}}
%     \Stilde(s_i),
%     \qquad
%     \hat{a}
%     =
%     a_{\hat{t}}.
%     \label{eq:prediction}
% \end{equation}



% \newpage
% Our pipeline produces, for every step $s_t$, a scalar score $\Stilde(s_t)$ whose argmax is the predicted decisive-error step (\eqref{eq:prediction}). The score is the output of \soap's rescoring phase (\S\ref{sec:method:recap}), which corrects an off-the-shelf per-step base score $\Sbase$ (\S\ref{sec:method:base}) by backpropagating each step's anomaly to its upstream sources along attention-derived causal-dependency weights $w_{i,t}$ (\S\ref{sec:method:weights}).

% \subsection{Per-step representations and base scoring}
% \label{sec:method:base}

% \textbf{Representation extraction.} For each step $s_t$, we form an input by concatenating the predecessor steps $s_1,\ldots,s_{t-1}$ with $s_t$ and run a single forward pass of $\proxy$. Let $\mathcal{T}_t$ denote the token positions occupied by $s_t$ and let $h^{(l)}_p \in \mathbb{R}^d$ be the hidden state at layer $l$ and position $p$. We mean-pool the hidden states over the step's own positions at a chosen layer $l^\star$:
% \begin{equation}
%   v_t \;=\; \operatorname*{mean}_{p \in \mathcal{T}_t}\; h^{(l^\star)}_p \;\in\; \mathbb{R}^d,
%   \label{eq:embedding}
% \end{equation}
% where the layer $l^\star$ is a hyperparameter (\S\ref{sec:experiments}).

% \textbf{SVD projection score.} We score $v_t$ by its squared projection onto a band of right singular vectors of a reference matrix $R\in\mathbb{R}^{N\times d}$ of step embeddings. Let $R = U\Sigma V^\top$ be the thin SVD, with columns of $V$ ordered by decreasing singular value. The directions that separate erroneous from non-erroneous steps lie in a contiguous band $\mathcal{C}=\{c_\text{begin},\ldots,c_\text{end}-1\}$ of singular vectors, selected as a hyperparameter.
% % Leading components capture coarse shared structure (prose style, agent voice, formatting) and trailing components capture noise; the informative band therefore sits in the mid-spectrum.
% We define the projection magnitude
% \begin{equation}
%   \pi(v_t) \;=\; \operatorname*{mean}_{c \in \mathcal{C}}\; \bigl\langle v_t,\, V_{:,c} \bigr\rangle^{2}.
%   \label{eq:proj}
% \end{equation}
% Erroneous steps yield \emph{smaller} projection magnitudes on the selected band, so we obtain the base score by passing $\pi$ through a monotonically decreasing orientation map $\rho$ (e.g., $\rho(x)=-x$):
% \begin{equation}
%   \Sbase(s_t) \;:=\; \rho\bigl(\pi(v_t)\bigr),
%   \label{eq:base-score}
% \end{equation}
% so that larger $\Sbase$ is more error-like, consistent with \eqref{eq:prediction}. A norm baseline replaces $\pi(v_t)$ with $\|v_t\|_2$; the rescoring below is agnostic to the form of $\Sbase$.

% \subsection{Attention-mass dependency weights}
% \label{sec:method:weights}

% \soap requires, for every ordered pair $(i,t)$ with $i<t$, a non-negative coefficient $w_{i,t}$ quantifying how strongly step $t$ depends on step $i$. We derive it from the proxy's post-softmax attention on the same per-step inputs. Let $A^{(l,h)}\in[0,1]^{N\times N}$ be the causally masked attention at layer $l$, head $h$, row-stochastic over keys. The attention mass routed from step $t$ into step $i$ is the mass a typical step-$t$ query places on step $i$'s positions:
% \begin{equation}
%   m^{(l,h)}_{i,t} \;=\; \operatorname*{mean}_{p \in \mathcal{T}_t}\; \sum_{q \in \mathcal{T}_i} A^{(l,h)}_{p,q}.
%   \label{eq:attn-mass-lh}
% \end{equation}
% % The asymmetry is deliberate: averaging over $\mathcal{T}_t$ gives the typical mass a step-$t$ query routes into step $i$, while summing over $\mathcal{T}_i$ retains the total probability landing in step $i$; dividing also by $|\mathcal{T}_i|$ would penalize long predecessors at fixed routing intensity, the opposite of the dependency we want.
% We average over heads and a selected layer band $L^\star$,
% \begin{equation}
%   m_{i,t} \;=\; \operatorname*{mean}_{l \in L^\star,\; h}\; m^{(l,h)}_{i,t},
%   \label{eq:attn-mass-agg}
% \end{equation}
% and renormalize over predecessors into a row-stochastic weight,
% \begin{equation}
%   w_{i,t} \;=\; \frac{m_{i,t}}{\sum_{j < t} m_{j,t}}.
%   \label{eq:weight}
% \end{equation}
% The renormalization absorbs any mass routed to non-predecessor positions (BOS, chat-template tokens, $s_t$'s own tokens). A softmax-sharpening variant $w_{i,t}=\mathrm{softmax}_j(m_{j,t}/\tau)_i$ is also supported (ablated in \S\ref{sec:experiments}).

% \subsection{Rescoring by causal attention backpropagation}
% \label{sec:method:recap}
% \textbf{Downstream contamination problem.} A per-step scorer evaluates each step $s_t$ together with the steps that precede it. Once the decisive error occurs at $t^\star$, every later step is generated from a context that already contains the mistake, so the scorer registers elevated anomaly not only at $t^\star$ but at many steps after it. Taking the argmax of $\Sbase$ therefore often returns a late step that merely inherits the corruption, rather than the step that introduced it. Figure~\ref{fig:offset} illustrates this on the hand-crafted subset of Who\&When: base-scorer predictions land systematically \emph{after} the annotated decisive step. \TODO{Use a motivating example instead}
% \textbf{Formulation.} Motivated by the downstream contamination phenomenon, we treat the contaminated suffix not as noise to be suppressed but as evidence to be returned to its source: each step's anomaly is backpropagated along the dependency weights to the upstream steps it was inherited from. Concretely, every step collects the anomaly of its dependents, weighted by how strongly they attended to it:
% \begin{equation}
%   \Stilde(s_i) \;=\; \Sbase(s_i) \;+\; \gamma\, \frac{\sum_{t>i} w_{i,t}\, \Sbase(s_t)}{\sum_{t>i} w_{i,t}}, \qquad \Stilde(s_T) = \Sbase(s_T),
%   \label{eq:recap}
% \end{equation}
% where $\gamma \in [0,1]$ controls the propagation strength. Setting $\gamma=0$ recovers the base scorer exactly; the final step, having no dependents, is left unchanged. The normalization by the total attention received, $\sum_{t>i} w_{i,t}$, keeps the correction honest: without it, any heavily attended step---a plan, a restatement of the task---would accumulate blame simply for being popular. Dividing by attention received turns the correction into the mean anomaly of a step's dependents per unit of attention, so a step is promoted only when its dependents are anomalous \emph{relative to how much they leaned on it}. The right-hand side reads the original $\Sbase$, never the rescored value, so the correction is single-pass and does not cascade through the trajectory. As a result, a decisive error whose own representation looks unremarkable is still promoted whenever the steps built on it are broken---which is precisely what makes an error decisive. The prediction follows from \eqref{eq:prediction}.

% \textbf{Properties.} \soap introduces a single new continuous hyperparameter $\gamma$ on top of the attention layer band $L^\star$. It requires no training data, no fine-tuning of $\proxy$, and no per-trajectory calibration: scores are ranked within each trajectory independently, sidestepping calibration across trajectories of differing length and domain.